\head{第十一版新核查：第0章21.1--21.4}
\source{EGA IV第一部分，分册印刷第153--164页／上传PDF第152--163页；底本为用户上传的法文扫描本}
本轮连续译入§21章首、（21.1.1）--（21.4.8），其中包括子编号引理（21.1.6.1）和（21.4.0），共34个原著编号条目。
21.5开头仅用于确认下一接续点，没有计入完成范围。

本轮从实际可读取的第十版工程接续，不以聊天中的完成声明代替源码。新增正文与源页逐段对照；
在可核实处将正确表述写入正文，原文差异置于脚注。下面的“校订”包含不同性质的问题：
原断言在字面假设下有反例、错误的中间等式、需展开的半线性记号、内部引用及定义约定。
\textbf{不能把所有表中项目相加，作为发现了同样数量的错误定理。}

\section*{处理清单与边界}
\addcontentsline{toc}{section}{处理清单与边界}
\begin{longtable}{@{}p{.17\linewidth}p{.17\linewidth}p{.57\linewidth}@{}}
\toprule 标识 & 位置 & 处理\\
\midrule\endhead
K11-E01 & 21.1.10，ii & 字面表述存在两个独立问题：集合并会丢失重合索引；中间扩张的$p$-基系数环不能任意扩大。给出两例及修订证明。\\
K11-T01 & 21.2.2 & 平凡扩张结论写为$A$-平凡，并指出原文末句的$B$-平凡记号问题。\\
K11-T02 & 21.2.5 & 微分模不变性的内部引文由21.1.5，iii改为ii。\\
K11-T03 & 21.3.6 & 将$j\otimes1$展开为真正相容的Frobenius扭曲公式；证明两个基变换同构保留。\\
K11-T04 & 21.4.0 & 准确表述$K^p[k]$在$K^p$上整，避免把未必包含$K^p$的$k$当作其扩域。\\
K11-T05 & 21.4.5 & 删除中间域的两个不同相对微分模之间的不当等式；给出反例并核验定理证明不受影响。\\
K11-G01--G04 & 对应条目 & 展开Hom分裂、指定$p$-基像的提升、修订引理的后续应用及Zorn上界步骤。\\
K11-N01--N06 & 对应条目 & 保留历史定义、Frobenius标量作用、幂零元理想与幂零理想、形式光滑性假设、正合列末端和可分性判据的区别。\\
\bottomrule
\end{longtable}
这些项目是本次译者分析，并非一份新取得的作者勘误。核对第四卷第四部分的《勘误与增补（清单三）》时，
须严格区分$0_{\mathrm{IV}}$的21节与第IV章的21节；后者有关除子的同号条目不是本轮微分代数条目的勘误。

\audit{K11-E01}{两级$p$-基与集合并：两个条件都不能省略}{按原文的字面假设，等价关系不成立；正文已写入校订形式}
\oldauditsrc{$0_{\mathrm{IV}}$（21.1.10），分册印刷第157页／上传PDF第156页}
原文设$B^p\subseteq C\subseteq B$，$C$为$A$-子代数，$M\subseteq C$、$N\subseteq B$。
其$(ii)$仅假设$M$为$C$在$A$上的$p$-基，便断言
\[
 N\text{在}C\text{上$p$-无关}
 \quad\Longleftrightarrow\quad
 M\cup N\text{在}A\text{上$p$-无关}.
\]
紧接着的证明写道可归结于$B^p\subseteq A\subseteq C\subseteq B$。
但从单个扩张的定义出发替换基环时，不能同时不加说明地保留另一个扩张的$p$-基。
此外，集合并会将$M\cap N$中的同一元素只计一次，与乘积基证明所使用的两个独立索引不符。

\lead{反例一：只检验重合元素问题。}
固定素数$p$和超越元$t$，令
\[
 A=\mathbf F_p(t^p),\qquad C=B=\mathbf F_p(t),\qquad M=N=\{t\}.
\]
此时$B^p=A$，故系数环问题已经不存在。
令$s=t^p$，则$X^p-s$在$\mathbf F_p(s)[X]$中不可约：
若$s$是该域中某个元素的$p$次幂，则它在$s=0$处的赋值应是$p$的倍数，但该赋值为$1$。
因此$[B:A]=p$，$1,t,\ldots,t^{p-1}$为$A$-基，$M$为$C$在$A$上的$p$-基。

集合并$M\cup N$仍为$\{t\}$，所以在$A$上$p$-无关。
但是$N$不是在$C$上$p$-无关，因为其受限单项式族包含$1,t$，且有非平凡$C$-线性关系
\[
 (-t)\cdot1+1\cdot t=0.
\]
这直接反驳原文的反向蕴含。

\lead{反例二：两集合不交，仍不能扩大中间系数环。}
取
\[
 A=\mathbf F_p,\qquad C=\mathbf F_p(t^p),\qquad B=\mathbf F_p(t),
 \qquad M=\{t^p\},\qquad N=\{t\}.
\]
有$B^p=C$，故满足原文$B^p\subseteq C$，且$M\cap N=\varnothing$。
$M$是$C$在$A$上的$p$-基，因为该定义中的系数环是
\[
 A[C^p]=\mathbf F_p(t^{p^2}),
\]
而$1,t^p,\ldots,t^{(p-1)p}$是$C$在这一域上的基。
同样，$N$是$B$在$C$上的$p$-基。

但检验$M\cup N$在$A$上的$p$-无关性时，实际系数环是
$A[B^p]=\mathbf F_p(t^p)$。其中已含元素$t^p$，于是
\[
 1\cdot(t^p)+(-t^p)\cdot1=0
\]
是一条非平凡的系数关系。所以即使两集合不交，原文的正向蕴含仍然失败。

\lead{正文采用的修订形式及证明。}
令$R=A[B^p]\subseteq C$。将$(ii)$改为：
若$M$为$C$在$R$上的$p$-基，且$M\cap N=\varnothing$，则原等价成立。
因为$C^p\subseteq B^p\subseteq R$，$M$的受限单项式族$P$是$C$的$R$-基；
因为$B^p\subseteq C$，$N$的受限单项式族$Q$在$C$上线性无关，恰好表达$N$的$p$-无关性。

若$Q$在$C$上线性无关，对有限关系
\[
 \sum_{\lambda,\mu}a_{\lambda\mu}x_\lambda y_\mu=0,
 \quad a_{\lambda\mu}\in R,\quad x_\lambda\in P,\quad y_\mu\in Q,
\]
先按$Q$收集系数，再按$P$展开，便得每个$a_{\lambda\mu}=0$。
反之，将$Q$的任何有限$C$-线性关系的系数按$P$展开，乘积族的独立性又推出原系数全零。
$M,N$不交保证乘积族与并集的受限单项式族按索引一一对应。

也可以不要求集合不交，而改用保留两套索引的\emph{拼接族}；但这不是原文的去重集合并，
因此正文选用直接明写条件的形式。
如果把原文前一段的“可假设$B^p\subseteq A$”解释为已经对本条施加的约定，
那么关于系数环的修订是在正文中明确该约定；仍不能据此消除反例一中的集合重合问题。

\lead{后续应用没有被一并否定。}
（21.3.5）用到的是塔$A^p\subseteq B^p\subseteq A$以及$B^p\subseteq A\subseteq B$。
这两次应用中，最大环的$p$次幂已在最小环中；两级相对$p$-基也不交。
（21.4.2）加入的元素则明确不属于先前生成的子域。因此这些后续应用均满足正文的修订条件。
原文$(i)$关于生成系的结论不需要两集合不交，且仍成立：
$C=A[C^p,M]=R[M]$、$B=C[N]=R[M,N]$。

\audit{K11-T03}{半线性$j$的基变换不能写成普通$j\otimes1$}{展开并校正记号；原来的典范同构保留}
\oldauditsrc{（21.3.6，ii），分册印刷第162页／上传PDF第161页}
原文有
\[
 A\xrightarrow{i}B\xrightarrow{j}A,\qquad ji=F_A,\quad ij=F_B,
\]
然后对$f:A\to A'$、$B'=B\otimes_A A'$直接写$j'=j\otimes1:B'\to A'$。
通常的张量记号要求相应的平衡关系相容；这里$j$相对于$A$是Frobenius半线性的，不能省去这一信息。

\lead{普通张量解释为何不成立。}
取$A=B=\mathbf F_p[t]$、$i=1$、$j=F_A$，并取$A'=A$、$f=1$。
在$B\otimes_A A'$中，$t\otimes1=1\otimes t$。
若按$b\otimes a'\mapsto j(b)a'$解释原式，两者分别映到$t^p$和$t$，并不相等。
因此这不是一个良定义的映射。

\lead{正确公式及全部相容性。}
令
\[
 i'(a')=1\otimes a',\qquad
 j'(b\otimes a')=f(j(b))(a')^p.
\]
先检验平衡关系。对$a\in A$，
\begin{align*}
 j'(i(a)b\otimes a')
 &=f(j(i(a)b))(a')^p\\
 &=f(a)^p f(j(b))(a')^p\\
 &=j'(b\otimes f(a)a').
\end{align*}
加法性来自特征$p$下的Frobenius加法性；乘法和单位相容性由公式直接计算。
故$j'$是环同态。又
\[
 j'i'(a')=(a')^p,
\]
以及
\begin{align*}
 i'j'(b\otimes a')
 &=1\otimes f(j(b))(a')^p\\
 &=i(j(b))\otimes(a')^p\\
 &=b^p\otimes(a')^p=(b\otimes a')^p.
\end{align*}
所以$j'i'=F_{A'}$、$i'j'=F_{B'}$。
若$g:B\to B'$是$b\mapsto b\otimes1$，还满足$j'g=fj$及$i'f=gi$，
恰好形成原文$(i)$所需的交换图。

\lead{为什么$\Theta$的基变换同构仍成立。}
写出张量所用的结构同态，便有
\begin{align*}
 \Theta_{B'/A'}
 &=\Omega^1_{B'/A'}\otimes_{B',j'}A'\\
 &\simeq(\Omega^1_{B/A}\otimes_{B,g}B')\otimes_{B',j'}A'\\
 &\simeq\Omega^1_{B/A}\otimes_{B,fj}A'\\
 &\simeq(\Omega^1_{B/A}\otimes_{B,j}A)\otimes_{A,f}A'\\
 &=\Theta_{B/A}\otimes_A A'.
\end{align*}
这正是（21.3.6.1），不需要平坦性假设。

\lead{Frobenius扩张标量的第二个同构。}
对任意$B$-模$M$，有
\[
 (M\otimes_{B,F_B}B)\otimes_B B'
 \simeq(M\otimes_B B')\otimes_{B',F_{B'}}B'.
\]
可将映射写成
\[
 (m\otimes b)\otimes c'\longmapsto(m\otimes1)\otimes g(b)c'.
\]
反向映射为
\[
 (m\otimes b')\otimes c'\longmapsto(m\otimes1)\otimes(b')^p c'.
\]
两者的平衡性由$gF_B=F_{B'}g$得到；逐个张量因子复合即为恒等。
取$M=\Omega^1_{B/A}$并结合微分模的基变换，即得（21.3.6.2）。
本项并不声称原文这两个同构错误；问题在于$j'$的普通张量写法没有表示出必需的扭曲。

\audit{K11-T05}{中间域上的两个相对微分模并不相同}{原证明中的等式不成立；定理结论不变}
\oldauditsrc{（21.4.5）证明首段，分册印刷第163页／上传PDF第162页}
原文令$K'=k(K^p,(x_\alpha))$，并在首段把
\[
 \Omega^1_{K'/k(K^p)}\quad\text{与}\quad\Omega^1_{K'/k}
\]
等同。对整个域$K$，$\Omega^1_{K/k}=\Omega^1_{K/k(K^p)}$确实成立，
因为$K^p$的元素在$K$中是$p$次幂；但对中间域$K'$，同样的理由不能直接使用。
$K^p$的元素未必是$K'$中元素的$p$次幂。

\lead{一个直接反例。}
取$k=\mathbf F_p$、$K=\mathbf F_p(t)$，取空元素族，则
$K'=K^p=\mathbf F_p(s)$，其中$s=t^p$。
于是
\[
 \Omega^1_{K'/k(K^p)}=\Omega^1_{K'/K'}=0,
 \qquad \Omega^1_{K'/k}=K'\,ds\ne0.
\]
$ds\ne0$可直接用$k$-导子$\partial/\partial s$验证，因为它将$s$映到$1$。
因此原等式不能保留在正文。

\lead{修订后证明仍闭合。}
实际只需知道$d_{K'/k(K^p)}(x_\alpha)$生成
$\Omega^1_{K'/k(K^p)}$，这一点由$K'$的生成方式得到。
若$d_{K/k}(x_\alpha)$生成$\Omega^1_{K/k}$，正合列
\[
 \Omega^1_{K'/k(K^p)}\otimes_{K'}K
 \longrightarrow\Omega^1_{K/k(K^p)}
 \longrightarrow\Omega^1_{K/K'}\longrightarrow0
\]
的左箭头满射，故$\Omega^1_{K/K'}=0$。
由（21.4.4），$K=K'(K^p)=K'$，从而得到$p$-生成性质。
独立性部分则按原证明利用$K_\alpha$及正合列，反证$d(x_\alpha)$不能落入其余微分生成的子空间。
这些论证不需要上述错误等式。因此正文删去等同关系，写入真正需要的相对微分模。

\audit{K11-N04}{有$p$-基不自动意味着形式光滑}{原定理条件提醒；不列为原文错误}
\oldauditsrc{（21.2.6）--（21.2.7），分册印刷第159--160页／上传PDF第158--159页}
令$A=\mathbf F_p$、$B=\mathbf F_p[t]/(t^p)$。
每个$p$次幂都是常数，所以$B^p=\mathbf F_p$，且
$1,t,\ldots,t^{p-1}$是$B$的$\mathbf F_p$-基。因此$\{t\}$确为$p$-基。
但（21.2.6.2）就是Frobenius映射$B\to B$，它把非零的$t$映为零，故不单射。

取
\[
 E=\mathbf F_p[T]/(T^{2p}),\qquad I=(T^p)\subset E,
 \qquad E/I\simeq B.
\]
$I^2=0$。若恒等映射$B\to E/I$有$A$-代数提升，$t$必须映到某个
$z=T+T^ph(T)$。特征$p$下
\[
 z^p=T^p+T^{p^2}h(T)^p=T^p\ne0\quad\text{于}E,
\]
因为$p^2\ge2p$；但$B$中$t^p=0$，矛盾。
这证明$B/A$并非形式光滑，同时具体显示了原定理条件$(i)$的作用。
正文没有删除该条件，也不把相对于$B^p$的形式光滑性与在$A$上的绝对提升条件混为一谈。

\audit{K11-N06}{微分模为零与可分性是不同的判据}{原文使用边界；不是新反例推翻原定理}
\oldauditsrc{（21.4.4），并对照已译（20.6.3）}
令
\[
 k=\mathbf F_p(t),\qquad
 K=\bigcup_{n\ge0}\mathbf F_p(t^{1/p^n}).
\]
对每个$K$中的元素，总能在后一级找到$p$次方根，因此$K^p=K$。
（21.4.4）给出$\Omega^1_{K/k}=0$。
然而$t^{1/p}\notin k$，其最小多项式是$X^p-t$，所以$K/k$是非平凡纯不可分代数扩张。

非完美性模的判据没有失效：$K$是完美域，$\Omega^1_K=0$，故
\[
 \Upsilon_{K/k}
 =\Ker(\Omega^1_k\otimes_k K\to\Omega^1_K)
 =K\,dt\ne0.
\]
所以（20.6.3）准确检出不可分性。不能把（21.4.4）的微分消失判据，替换为任意域扩张的可分性判据。

\audit{K11-T01/T02/T04}{三处局部表述和引文的校订}{没有合计为三个错误定理}
\oldauditsrc{（21.2.2）、（21.2.5）、（21.4.0）}
\lead{21.2.2的平凡扩张。}
给定的是$B$的$A$-扩张$E$，并假设在$B^p$上平凡。
（21.2.1）的结论是存在$A$-代数截面$B\to E$，所以所需结论是$A$-平凡。
原文末句的“$B$-平凡”不应使读者误以为已另有给定的$B$-代数结构。

\lead{21.2.5的内部引用。}
原扫描页引用（21.1.5，iii），但以$A[B^p]$替换$A$而保持微分模的性质在$(ii)$；
$(iii)$讨论的是乘法核的幂零元性质。正文采用$(ii)$。

\lead{21.4.0的复合域。}
$k[K^p]=k(K^p)$的结论正确，但原文理由“$k$在$K^p$上代数”应准确化。
例如$k=\mathbf F_p(s)$、$K=\mathbf F_p(s,t)$，则$K^p=\mathbf F_p(s^p,t^p)$并不包含于$k$。
实际有意义的是$K^p[k]$在$K^p$上整；每个非零$z\in K^p[k]$满足
\[
 z^{-1}=\frac{z^{p-1}}{z^p}\in K^p[k],
\]
所以该环为域。这保留原结论而避免不合法的域扩张记号。

\audit{K11-G01/G02/G04}{补明三个证明步骤}{展开原有论证，不代写为原作者的新结论}
\lead{Hom满射与分裂。}
在（21.2.4）中，记$U=\Omega^1_{A/B^p}\otimes_A B$、$V=\Omega^1_{B/B^p}$，
$\theta:U\to V$为自然映射。（21.2.3）说明对任意$B$-模$L$，
$\Hom_B(V,L)\to\Hom_B(U,L)$满射。取$L=U$，则$1_U$延拓为$r:V\to U$，
满足$r\theta=1_U$。这直接证明$\theta$单射且分裂。
另由$\Der_A(B,L)\simeq L^I$，表示性给出
$\Omega^1_{B/A}\simeq\bigoplus_I B$，其基是对应的$dx_\alpha$。

\lead{指定$p$-基像的相容提升。}
平方零情形中，任何$A$-代数截面$w$对$p$次幂的值都已确定，因而自动与正文构造的$A[B^p]$-结构相容。
对幂零理想，逐次提升到$E/I^r$，每次指定$x_\alpha$的像为固定$z_\alpha$的像，
便得到每一步唯一的提升。对一般可容许环，离散商上的同态因这种唯一性自动与过渡同态相容，
从而可以取逆极限。这里没有从“各层存在提升”直接跳到“相容提升存在”。

\lead{Zorn引理的上界。}
对按包含排序的$p$-无关子集链，链的并仍$p$-无关：一个有限线性关系仅涉及有限多个元素，
它们都包含在链的某一个成员中，于是关系只能平凡。这给出（21.4.2）所需的上界。

\section*{定义约定及外部交叉核验}
\addcontentsline{toc}{section}{定义约定及外部交叉核验}
（21.1.1）--（21.1.2）在$p=0$时要求存在$\mathbf Q\to A$，比通常“$1_A$的加法阶无限”的条件强。
正文明确这是本节约定；没有将历史定义误判为错误。
（21.3.4）的序列右端没有$0$，因此不从“正合”臆造最右箭头的满射性。

外部仅以The Stacks Project的\href{https://stacks.math.columbia.edu/tag/07P1}{Tag 07P1}
和\href{https://stacks.math.columbia.edu/tag/07P2}{Tag 07P2}核对域扩张中$p$-基的定义、存在性及其与微分基的等价。
这些标准结论与本轮21.4的主线一致；不将外部材料当作本轮具体错误修订的作者依据，也不替换法文底本。
查询日期：2026年11月15日。

\section*{本轮技术检查的含义}
\addcontentsline{toc}{section}{本轮技术检查的含义}
源条目、公式标签、实际构建页数和日志检查记录保存在工程的\texttt{checks}目录。
编译成功只检验TeX和排版，不证明数学命题。
有限特征下的计算脚本只核验指定例子的多项式恒等式，完整一般论证以本报告给出的证明为准。
本轮没有重新逐句审校旧版全部正文和历史报告，也没有宣称全部上传材料已经译完。
